\documentclass[11pt]{article}
% ===========================================================================
%  Shared preamble for the AP Calculus AB exam-preparation worksheets.
%  Input right after \documentclass. Files beginning with "_" are not compiled
%  on their own. Navy and gold, to match the course's lesson pages.
% ===========================================================================
\usepackage[letterpaper,top=2.3cm,bottom=2.2cm,left=1.7cm,right=1.7cm,headheight=28pt,footskip=30pt]{geometry}
\usepackage{amsmath,amssymb}
\usepackage[T1]{fontenc}
\usepackage[utf8]{inputenc}
\usepackage{lmodern}
\usepackage{xcolor}
\usepackage[most]{tcolorbox}
\usepackage{enumitem}
\usepackage{fancyhdr}
\usepackage{array}
\usepackage{tikz}
\usepackage{pgfplots}
\pgfplotsset{compat=1.18}

\definecolor{navy}{HTML}{161E54}
\definecolor{ink}{HTML}{27357E}
\definecolor{gold}{HTML}{B98A00}
\definecolor{goldbg}{HTML}{FFFDF3}
\definecolor{inkbg}{HTML}{F6F7FD}
\definecolor{rule}{HTML}{C9CFF0}
\definecolor{okgreen}{HTML}{1B7A44}
\definecolor{okbg}{HTML}{F1F8F3}

\pagestyle{fancy}
\fancyhf{}
\fancyhead[L]{\footnotesize Name: \underline{\hspace{4cm}}}
\fancyhead[C]{\textbf{\color{navy}Integration Academy}\\[-2pt]{\scriptsize Calculus AB \textperiodcentered{} AP\textsuperscript{\textregistered} Exam Preparation}}
\fancyfoot[L]{\tiny\color{black!55}\parbox[t]{0.78\textwidth}{AP\textsuperscript{\textregistered} is a registered trademark of the College Board, which was not involved in the production of, and does not endorse, this material. \textcopyright{} Integration Academy.}}
\fancyfoot[R]{\scriptsize\thepage}
\renewcommand{\headrulewidth}{1.2pt}
\renewcommand{\headrule}{\hbox to\headwidth{\color{gold}\leaders\hrule height \headrulewidth\hfill}}
\renewcommand{\footrulewidth}{0pt}
\setlength{\parindent}{0pt}
\setlength{\parskip}{4pt}

% Title block: \wstitle{code}{title}{unit line}
\newcommand{\wstitle}[3]{\begin{center}{\small\bfseries\color{gold}WORKSHEET #1}\\[3pt]{\LARGE\bfseries\color{navy}#2}\\[3pt]{\small\color{black!60}#3}\end{center}\vspace{2pt}}
% Section heading
\newcommand{\sech}[1]{\par\vspace{10pt}{\large\bfseries\color{navy}#1}\par\nobreak\vspace{2pt}{\color{rule}\hrule height0.8pt}\vspace{7pt}}
% A level tag printed beside a group of questions
\newcommand{\level}[1]{\par\vspace{6pt}{\footnotesize\bfseries\color{gold}\MakeUppercase{#1}}\par\nobreak\vspace{2pt}}

\newtcolorbox{ideabox}{colframe=ink,colback=inkbg,coltitle=white,colbacktitle=ink,fonttitle=\bfseries\small,title={Key ideas},arc=2pt,boxrule=0.8pt,left=8pt,right=8pt}
\newtcolorbox{example}[1]{colframe=rule,colback=white,coltitle=navy,colbacktitle=inkbg,fonttitle=\bfseries\small,title={#1},arc=2pt,boxrule=0.8pt,borderline west={3pt}{0pt}{ink},left=10pt,right=8pt}
\newtcolorbox{frbox}[1]{colframe=gold,colback=goldbg,coltitle=white,colbacktitle=gold,fonttitle=\bfseries\small,title={#1},arc=2pt,boxrule=0.8pt,breakable,left=8pt,right=8pt}
\newtcolorbox{solbox}{colframe=okgreen,colback=okbg,coltitle=white,colbacktitle=okgreen,fonttitle=\bfseries\small,title={Worked solutions},arc=2pt,boxrule=0.8pt,breakable,left=8pt,right=8pt}
\newtcolorbox{scorebox}[1]{colframe=navy,colback=white,coltitle=white,colbacktitle=navy,fonttitle=\bfseries\small,title={#1},arc=2pt,boxrule=0.8pt,breakable,left=8pt,right=8pt}

\newcommand{\soln}{\par\smallskip\textbf{\color{ink}Solution.}\ }
\newcommand{\result}[1]{\par\smallskip\textbf{#1}}
% blank working space of a given height
\newcommand{\work}[1]{\par\vspace{#1}}
% a point value in the margin of a scoring guide
\newcommand{\pts}[1]{\hfill{\footnotesize\bfseries\color{gold}[#1]}}

% Coordinate grid for graph questions: \begin{apgraph}{xmin}{xmax}{ymin}{ymax} ... \end{apgraph}
\newenvironment{apgraph}[4]{\begin{center}\begin{tikzpicture}\begin{axis}[
  axis lines=middle, xmin=#1, xmax=#2, ymin=#3, ymax=#4,
  grid=both, grid style={black!12}, xtick distance=1, ytick distance=1,
  ticklabel style={font=\scriptsize}, xlabel={$x$}, ylabel={$y$},
  width=10cm, height=7.2cm, samples=80, clip=false,
  every axis plot/.append style={line width=1.1pt, color=ink}]}{\end{axis}\end{tikzpicture}\end{center}}
\newcommand{\closeddot}[2]{\addplot[only marks, mark=*, mark size=2.2pt, color=ink] coordinates {(#1,#2)};}
\newcommand{\opendot}[2]{\addplot[only marks, mark=*, mark size=2.2pt, color=ink, mark options={fill=white}] coordinates {(#1,#2)};}

% Multiple-choice options: four across, or stacked when the options are long.
\newcommand{\opts}[4]{\par\vspace{3pt}\begin{tabular}{@{}p{0.235\linewidth}p{0.235\linewidth}p{0.235\linewidth}p{0.235\linewidth}@{}}(A)\ #1 & (B)\ #2 & (C)\ #3 & (D)\ #4\end{tabular}\par\vspace{9pt}}
\newcommand{\optsv}[4]{\par\vspace{3pt}\begin{tabular}{@{}l@{\ \ }p{0.9\linewidth}@{}}(A) & #1\\(B) & #2\\(C) & #3\\(D) & #4\end{tabular}\par\vspace{9pt}}

\fancyhead[R]{\footnotesize\color{black!60}Worksheet 1.1}
\begin{document}

\wstitle{1.1}{The Instant Problem: From Average Rates to Limits}{Unit 1: Limits and Continuity}

\begin{ideabox}
\begin{itemize}[leftmargin=*,itemsep=2pt,topsep=2pt]
  \item The \textbf{average rate of change} of $f$ on $[a,b]$ is $\dfrac{f(b)-f(a)}{b-a}$, with units of $f$ per unit of $x$.
  \item An \textbf{instantaneous} rate is the limit of average rates as the interval shrinks to a point.
  \item $\displaystyle\lim_{x\to c}f(x)$ describes the values of $f$ \emph{near} $c$. The value $f(c)$ plays no part.
  \item A two-sided limit exists exactly when the left-hand and right-hand limits exist and are equal.
\end{itemize}
\end{ideabox}

\sech{Worked examples}
\begin{example}{Example 1 \textperiodcentered{} Average rate from a formula}
Find the average rate of change of $f(x)=x^2-5x$ on $[2,6]$.\soln $f(6)=6$ and $f(2)=-6$, so $\dfrac{6-(-6)}{6-2}=\dfrac{12}{4}$.\result{The average rate of change is $3$.}
\end{example}

\begin{example}{Example 2 \textperiodcentered{} Shrinking the interval}
For $s(t)=2t^2+t$, simplify the average rate of change on $[3,3+h]$, $h\neq0$, and find what it approaches as $h\to0$.\soln $\dfrac{2(3+h)^2+(3+h)-21}{h}=\dfrac{13h+2h^2}{h}=13+2h$.\result{As $h\to0$ the average rates approach $13$.}
\end{example}

\begin{example}{Example 3 \textperiodcentered{} Interpreting in context}
The depth of water in a pool, $D(t)$ cm, is measured at times $t$ minutes.\begin{center}\renewcommand{\arraystretch}{1.25}\begin{tabular}{|c|c|c|c|c|}\hline $t$ & $0$ & $4$ & $10$ & $15$ \\ \hline $D(t)$ & $60$ & $52$ & $34$ & $29$ \\ \hline\end{tabular}\end{center}Find the average rate of change of $D$ on $[4,10]$ and interpret it.\soln $\dfrac{D(10)-D(4)}{10-4}=\dfrac{34-52}{6}=-3$.\result{From $t=4$ to $t=10$ minutes, the depth of the water decreased at an average rate of 3 cm per minute.}
\end{example}

\newpage
\sech{Practice}

\level{Core}
\begin{enumerate}[leftmargin=*,itemsep=2pt]
  \item Find the average rate of change of $f(x)=x^3$ on $[-1,2]$.\work{2.3cm}
  \item Find the average rate of change of $f(x)=\dfrac1x$ on $[1,4]$.\work{2.3cm}
  \item A function has $\displaystyle\lim_{x\to 3^-}f(x)=5$, $\displaystyle\lim_{x\to 3^+}f(x)=5$ and $f(3)=1$. State $\displaystyle\lim_{x\to 3}f(x)$.\work{2.3cm}
  \item A function has $\displaystyle\lim_{x\to 0^-}g(x)=-2$ and $\displaystyle\lim_{x\to 0^+}g(x)=2$. Does $\displaystyle\lim_{x\to 0}g(x)$ exist? Explain.\work{2.3cm}
\end{enumerate}

\level{Exam level}
\begin{enumerate}[leftmargin=*,itemsep=2pt,start=5]
  \item For $s(t)=t^2-3t$, simplify the average rate of change on $[2,2+h]$, $h\neq0$, and state the instantaneous rate of change at $t=2$.\work{2.8cm}
  \item Let $g(x)=\dfrac{x^2-1}{x-1}$ for $x\neq1$ and $g(1)=7$. Find $\displaystyle\lim_{x\to 1}g(x)$.\work{2.8cm}
  \item Using the pool table in Example 3, find the average rate of change of $D$ on $[10,15]$ and interpret it with units.\work{2.8cm}
  \item The average rate of change of $f$ on $[2,5]$ is $-4$, and $f(5)=1$. Find $f(2)$.\work{2.8cm}
\end{enumerate}

\begin{frbox}{Question 9 \textperiodcentered{} Free-response practice}
The volume of water in a tank, $V(t)$ litres, is recorded at selected times $t$ minutes.\begin{center}\renewcommand{\arraystretch}{1.25}\begin{tabular}{|c|c|c|c|c|}\hline $t$ & $0$ & $3$ & $8$ & $12$ \\ \hline $V(t)$ & $200$ & $170$ & $130$ & $122$ \\ \hline\end{tabular}\end{center}
\begin{enumerate}[label=(\alph*),leftmargin=*,itemsep=2pt]
  \item Find the average rate of change of $V$ on $[3,8]$. Indicate units of measure.\work{2.6cm}
  \item Using correct units, interpret your answer to part (a) in context.\work{2.6cm}
  \item On which of the intervals $[0,3]$, $[3,8]$, $[8,12]$ is the water draining fastest on average? Show the work that leads to your answer.\work{2.6cm}
\end{enumerate}
\end{frbox}

\newpage
\begin{solbox}
\begin{enumerate}[leftmargin=*,itemsep=4pt]
  \item $\dfrac{8-(-1)}{2-(-1)}=\dfrac93=3$.
  \item $\dfrac{\frac14-1}{4-1}=\dfrac{-\frac34}{3}=-\dfrac14$.
  \item $5$. The one-sided limits agree; the value $f(3)=1$ is irrelevant.
  \item No. The one-sided limits exist but are not equal.
  \item $\dfrac{(2+h)^2-3(2+h)-(-2)}{h}=\dfrac{h+h^2}{h}=1+h$, which approaches $1$.
  \item For $x\neq1$, $g(x)=x+1$, so the limit is $2$ (not $g(1)=7$).
  \item $\dfrac{29-34}{15-10}=-1$. From $t=10$ to $t=15$ minutes the depth decreased at an average rate of 1 cm per minute.
  \item $\dfrac{1-f(2)}{3}=-4$ gives $f(2)=13$.
  \item \begin{enumerate}[label=(\alph*),leftmargin=*,itemsep=3pt]
    \item $\dfrac{V(8)-V(3)}{8-3}=\dfrac{130-170}{5}=-8$ litres per minute.
    \item From $t=3$ to $t=8$ minutes, the volume of water in the tank decreased at an average rate of 8 litres per minute.
    \item $[0,3]$: $\dfrac{170-200}{3}=-10$. $[3,8]$: $-8$. $[8,12]$: $\dfrac{122-130}{4}=-2$. The water drains fastest on $[0,3]$, at an average of 10 litres per minute.
  \end{enumerate}
\end{enumerate}
\end{solbox}

\end{document}
